\section{The Syllables of the Gods}

I want to call attention to an aspect of Tomberg's work that is perhaps unnoticed. He writes so straightforwardly, so poetically, so lucidly, that one can obtain an illusion (particularly if you are only reading with your discursive intellect) that you are understanding him without, in fact, catching on. I know because I did this.

\begin{quotex}
And then I have cited a verse from the Song of the Soul of St John of the Cross, because it has the virtue of awakening the deeper layers of the soul...
\end{quotex}


There it is, at the beginning of the Magician, a few sentences into his work. He means this quite literally. He who wishes to awaken deeper layers of his soul can memorize, meditate upon, and at times recite this verse. In modern parlance, this poetry is a \textbf{\emph{spell}}. This is not to be taken in a mechanical way, as if one could (in the sleep state) merely mouth these words and expect great things to occur in one's interior life. However, it is quite literally true nonetheless. I suppose there are other verses of poetry which could do the same thing, such as Novalis' \emph{Hymns to the Night}, or passages from Wordsworth, which taught and inspired a large portion of the generation of Victorians, people whose souls were arguably a lot thicker and deeper than ours tend to be today.

By the very act and virtue (read this word as immanent/inherent potency or power, not merely sentimental moral quality) of having a sacred Latin language in which the Scriptures could be memorized during the Middle Ages, our spiritual ancestors imparted a species of magical and incantatory power to their thoughts, words, and meditations on the Divine. He who committed all of the Psalms to memory (as those who began their course of monastic study at Cassiodorus' school had to do, a precedent which held for successor schools) had training as a doctor, an exorcist, a priest, a sage, and perhaps even a magician.

This doesn't mean that we have to advocate a revival of Latin. Think of reading antique poetry or Scripture as an exercise in active-passive receptivity, as does the high priestess, who covers the intermediate realms with a veil, who is seated, who is reading a book, and who wears the tiara with three levels. She is not ``busy'' inventing narratives, exploring fantasies, or hammering out rationalizations/self-justifications. In other words, she is distinctly un-modern (in the shallow sense). We are perfectly cognizant of Goethe's famous dictum that a man who is not of his own times, cannot be a man of any, but what we are insisting on is a genuine multiculturalism, rather than one in name only. Is it not right to place ourselves in the shoes of our ancestors? To see the world from their eyes? CS Lewis got this far, but we go further\footnote{\url{http://www.gornahoor.net/?p=3545}}: we claim that in this act of transcending our own times, of traveling into the medieval and Renaissance world which gave birth to the French hermetic tradition, we will encounter the timelessness which is always \emph{there}.


Here is another verse of poetry that contains a magical motive power:

\begin{quotex}
\begin{verse}
A painter of the Umbrian school\\
designed upon a gesso ground\\
the nimbus of the baptized God:\\
the wilderness is cracked and brown.\\
But through the water pale and thin\\
still shine the unoffending feet\\
and there above the Painter set\\
the Father and the Paraclete.
\end{verse}
\end{quotex}


What does poetry set us down in the middle of? What is it that does not change, which is beyond even the syllables of the Gods\footnote{\url{http://www.amazon.com/Literature-Gods-Roberto-Calasso/dp/0375725431}}? Even if we say that Change is eternal and supreme, then this is what does not change. As Calasso writes, \emph{men always tell themselves, do as the gods do}. The Christian tradition regards that which is deathless, which changes not, as God. There are other names for it, including the Tao. \emph{That which suffers hell, and abides, is deathless. That which aspires to God, and is not shaken, is deathless} (my paraphrasing of Ur von Balthasar in \emph{Theo-Logic}). In poetry, we enter a word of multiplicity and duality, but without a loss of the living Logos, since within the syllables remains a path which draws us toward a remembered unity. Poetry (therefore) has been considered an antidote for the flux or dialectic of Reality that spirals downward: poetry attempts to reverse this flow and raise the reader or poet upward. Philosophy and religion (exotericism) make what is implicit in poetry more explicit, at the cost (sometimes) of vividness or innate vitality, but with a gain in discipline or form. The Autonomous Duality is challenged head-on.

Bonaventura says \emph{non coerci maximo continer tamen a minimo divinum est} (that which is not compelled by the most powerful, but which nevertheless can contain even the smallest thing, that is divine). Wagner has Hans Sach\footnote{\url{http://classicalmusic.about.com/od/classicalmusictips/qt/Was-Duftet-Doch-Der-Flieder-Lyrics-And-Text-Translation.htm}} sing,

\begin{quotex}
\begin{verse}
I feel it and can not understand it,\\
can keep it not, but not forgotten:\\
and I grasp it all, but I can measure it not!\\
But how should I take also\\
what seemed to me immeasurable?\\
No rule would fit there,\\
and yet no error did it contain.\\
It sounded so old -- and yet was so new\\
A birdsong as sweet as May!
\end{verse}
\end{quotex}


The fact that we ``discover'' so many things today is actually a sign of our stupidity. We have a long Fall to recover from. The fact that we have more knowledge of ancient times than our Medieval Age ancestors did doesn't help us appreciate or understand it any better. Actually, to endure a ``Dark Age'' and to follow the burning in one's heart (which is exactly the inner meaning of the Gothic spirituality that flowered in union with the ``new religion'') is a sign not of weakness, but of strength. If Time is flowing backwards from the Cross, then the Middle Ages were the restoration of Rome which had fallen in judgement for the sins it committed and the excesses it suffered: the Holy Roman Empire. What came before Rome? The Golden-Silvern-Brazen ages of Empire, and indeed one sees the restoration of these periods in the rise of Euro-Atlantis \& the rebirth of Russia, the persistence of China. And what was before that? More cataclysms, and even greater civilizations, including Solomon, Egypt, and Atlantis. The world is inundated with the Great Flood, under the sign of the Fish. Those who breath underwater will survive. God is re-creating the world.

Why is this? Is this not the triumph of the ``last man'', the craveling, and the weakling? We assert that water is the element of Life. This is why Lao-Tzu compares the Tao with water, because it goes to the lowest spot and generates life from below. It does not contend, and therefore, it wins the victory. The element of water is not mixed (as is earth), it is not perishable with what it consumes (as is fire), it is not evanescent (like the air). It is therefore, the first sign of life, of baptism, and of the upward spiral which (when it reflects what is above faithfully) partakes of the upward movement. Vicktor Shauberger\footnote{\url{https://www.youtube.com/watch?v=3SNnHyaLWi8}} has some interesting insights on water, how it generates power/life from inside of the spiral.

The High Priestess reminds us that our ``feminine'' side is not female, but rather active-passive, generating through creative imagination the needed ascent as it reflects in a virginal pool what comes from above. The Dyad\footnote{\url{http://www.gornahoor.net/?p=5557}} can be identified with the Wisdom (Sophia) of God. Together with the Monad, the Dyad is the parent of all which comes after.

Since all of us recapitulate the Macro-Cosm, it is necessary to find the virginal pool within ourselves. This is an act, as Plato taught, of remembering. I find what is already there, what is changeless, relative to the flow of thoughts, feelings, sensations which are multiple and are not ``I'', but rather than which the I perceives.

It is important not to mistake the congerie of little ``i''s or a particular configuration of these ``i''s (which are legion) for the real I. As Tomberg puts it, there are illusions ``above'' just as there are illusions ``below''. The more deeply and truly one strips away these mechanical programs (which operate automatically and easily and comfortably, not to mention deceptively), the firmer ground one is able to stand and then sit upon.

In the movement away from the ``I'' to the narrative of the ``we'', of liberation and of activism, to the time-bound provinciality of the modern ``hermetic'' who is sealed within a plastic and airtight capsule of great artifice, I see a parody of the Divine process which was intended to reawaken the soul. In this substitution is not Life, but death, because it contradicts the very highest spiritual laws of liberty, which is that one must first dare to the Unity, then listen to it properly.

Tomberg has already outlined the katalogical movement of the soul as it mimics the Tao of God, and catches fire with Truth. It is aspiration/mysticism, gnosis, magic, revelation. These are the syllables of the God, Yod-He-Vau-He.

Even dialectic itself has its place, but not as a substitute God or a false self.


\flrightit{Posted on 2014-05-17 by Logres}

